\documentclass[../../../include/open-logic-section]{subfiles}

\begin{document}

\olfileid{sth}{ordinals}{opps}
\olsection{Ordinal Penerus lan Ordinal Limit}

Ing sawetara bab sabanjure, ordinal bakal asring
kita gunakake. Mula migunani yen kita ngenalake
sawetara gagasan prasaja.

\begin{defn}
Kanggo sembarang ordinal $\alpha$, \emph{peneruse}
yaiku $\ordsucc{\alpha} =\alpha \cup \{\alpha\}$.
Kita ngarani $\alpha$ \emph{ordinal penerus} yen
$\ordsucc{\beta} = \alpha$ kanggo sawenehing
ordinal~$\beta$. Kita ngarani $\alpha$ \emph{ordinal
limit} yen lan mung yen $\alpha$ dudu ordinal kosong
lan dudu ordinal penerus.
\end{defn}
\noindent
Asil sabanjure nuduhake yen iki gagasan \emph{penerus}
kang trep:

\begin{prop}
Kanggo sembarang ordinal $\alpha$:
\begin{enumerate}
	\item $\alpha \in \ordsucc{\alpha}$;
	\item $\ordsucc{\alpha}$ iku ordinal;
	\item ora ana ordinal $\beta$ kang nyukupi
	$\alpha \in \beta \in \ordsucc{\alpha}$.
	\end{enumerate}
\end{prop}

\begin{proof}
Cetha yen $\alpha \in \alpha \cup \{\alpha\}
= \ordsucc{\alpha}$. Semono uga,
$\ordsucc{\alpha}$ iku himpunan transitif kang
anggotane ordinal, mula dadi ordinal miturut
\olref[basic]{corordtransitiveord}. Ora mungkin
$\alpha \in \beta \in \ordsucc{\alpha}$,
amarga saka kene $\beta \in \alpha$ utawa
$\beta = \alpha$, nentang \olref[basic]{ordordered}.
\end{proof}
\noindent
Iki uga ngidini wujud bukti induksi transfinit liyane:

\begin{thm}[Induksi Transfinit Prasaja]\ollabel{simpletransrecursion}
Upamane $\phi(x)$ rumus kang nyukupi:
\begin{enumerate}
	\item $\phi(\emptyset)$; lan
	\item kanggo sembarang ordinal $\alpha$, yen
	$\phi(\alpha)$ mula $\phi(\ordsucc{\alpha})$; lan
	\item yen $\alpha$ ordinal limit lan $(\forall \beta \in
	\alpha)\phi(\beta)$, mula $\phi(\alpha)$.
\end{enumerate}
Mula $\forall \alpha \phi(\alpha)$.
\end{thm}

\begin{proof}
Kita mbuktekake kontrapositife. Upamane ana ordinal
kang nyukupi $\lnot\phi$; pilih $\gamma$ minangka
ordinal paling cilik kang mangkono. Mula utawa
$\gamma = \emptyset$, utawa
$\gamma = \ordsucc{\alpha}$ kanggo sawijining
$\alpha$ kang nyukupi $\phi(\alpha)$, utawa
$\gamma$ ordinal limit lan
$(\forall \beta \in \gamma)\phi(\beta)$.
Saben kemungkinan iki nentang salah siji syarat,
mula ora ana ordinal kang luput saka rumus kasebut.
\end{proof}
\noindent
Siji notasi pungkasan iki bakal migunani mengko:

\begin{defn}\ollabel{defsupstrict}
Yen $X$ himpunan ordinal, mula
$\supstrict(X) = \bigcup_{\alpha \in X}
\ordsucc{\alpha}$.
\end{defn}
\noindent
Ing kene, ``lsub'' tegese ``wates ndhuwur ketat kang
paling cilik''.\footnote{Sawetara buku nganggo
``$\text{sup}(X)$'' kanggo iki. Nanging buku liyane
nganggo ``$\text{sup}(X)$'' kanggo wates ndhuwur
\emph{ora ketat} kang paling cilik, yaiku mung
$\bigcup X$. Yen $X$ nduweni anggota paling gedhe,
$\alpha$, loro gagasan iki beda: wates ndhuwur
\emph{ketat} kang paling cilik yaiku
$\ordsucc{\alpha}$, dene wates ndhuwur
\emph{ora ketat} kang paling cilik mung $\alpha$.}
Asil sabanjure njlentrehake iki:

\begin{prop}
Yen $X$ himpunan ordinal, $\supstrict(X)$ iku
ordinal paling cilik kang luwih gedhe tinimbang
saben ordinal ing $X$.
\end{prop}

\begin{proof}
Upamane $Y = \Setabs{\ordsucc{\alpha}}{\alpha \in X}$,
dadi $\supstrict(X) = \bigcup Y$. Amarga ordinal
transitif lan saben anggota ordinal uga ordinal,
$\supstrict(X)$ iku himpunan transitif kang anggotane
ordinal, mula ordinal miturut
\olref[basic]{corordtransitiveord}.

Yen $\alpha \in X$, mula $\ordsucc{\alpha} \in Y$,
dadi $\ordsucc{\alpha} \subseteq \bigcup Y
= \supstrict(X)$, saengga
$\alpha \in \supstrict(X)$. Mula $\supstrict(X)$
luwih gedhe kanthi ketat tinimbang saben ordinal
ing $X$.

Kosok baline, yen $\alpha \in \supstrict(X)$,
mula $\alpha \in \ordsucc{\beta} \in Y$ kanggo
sawenehing $\beta \in X$, dadi $\alpha \leq
\beta \in X$. Saben wates ndhuwur ketat kanggo
himpunan mau mesthi ngemot anggota-anggota iki, mula
$\supstrict(X)$ iku wates ndhuwur ketat kang
\emph{paling cilik} tumrap $X$.
\end{proof}

\end{document}
